\documentclass{../zkdl-presentation-template}

\title[Криптографія]{\textbf{Практичне №6.Асиметрична криптографія}}

\date{29.10.2026}

% --- Додатково для цієї презентації ---
\usetikzlibrary{calc,arrows.meta,shapes.misc}

\newcommand{\legendre}[2]{\left(\frac{#1}{#2}\right)}
\newcommand{\QR}{\mathcal{Q}}
\newcommand{\QNR}{\overline{\mathcal{Q}}}
\newcommand{\PSQ}{\widetilde{\mathcal{Q}}}
\newcommand{\inR}{\in_R}

% кольори для схем протоколів
\colorlet{alicecol}{oc-indigo-8}
\colorlet{bobcol}{oc-teal-8}
\colorlet{evecol}{oc-red-8}

% палітра "ghoul edition"
\definecolor{ghoulred}{HTML}{C0001E}
\definecolor{ghoulblood}{HTML}{6E0010}
\definecolor{ghoulbg}{HTML}{0A0A0C}
\definecolor{ghoulgray}{HTML}{8A8A90}

% слайд "ghoul edition": весь слайд малюється як фон у \ghoulbackground
% (визначається безпосередньо перед відповідним слайдом), \begin{frame}[plain, ghoul]
\usetikzlibrary{fadings}
\newcommand{\ghoulbackground}{}
% спільна основа ghoul-слайдів: фон, затемнення, заголовок (#1)
\newcommand{\ghoulbase}[1]{%
    \path[fill overzoom image=6-crypto-practice/background.png] (0,0) rectangle (\paperwidth,\paperheight);
    \fill[black, opacity=0.35] (0,0) rectangle (\paperwidth,\paperheight);
    \fill[black, opacity=0.75, path fading=north] (0,0) rectangle (\paperwidth,0.55\paperheight);
    \fill[black, opacity=0.6, path fading=south] (0,0.82\paperheight) rectangle (\paperwidth,\paperheight);
    \node[anchor=north west, font=\Large\bfseries, text=ghoulred!60!black] at (0.57,9.17) {#1};
    \node[anchor=north west, font=\Large\bfseries, text=white] (gtitle) at (0.50,9.22) {#1};
    \draw[ghoulred, line width=1.4pt] (0.50,8.5) -- (gtitle.east |- 0,8.5);
}
\tikzset{
    gtxt/.style={text=white, font=\small},
    glass/.style={fill=black, fill opacity=0.6, text opacity=1, rounded corners=2pt, inner sep=3pt},
    gmsg/.style={-{Stealth[length=2.8mm]}, line width=1.1pt, ghoulred}
}
% картка персонажа: #1 -- центр, #2 -- ширина, #3 -- файл
\newcommand{\ghoulcard}[3]{%
  \node[inner sep=0pt] (card) at #1 {\includegraphics[width=#2]{#3}};
  \draw[ghoulred, line width=2.2pt, opacity=0.35] (card.south west) rectangle (card.north east);
  \draw[ghoulred, line width=0.8pt] (card.south west) rectangle (card.north east);
}
\makeatletter
\define@key{beamerframe}{ghoul}[true]{\setbeamertemplate{background}{\ghoulbackground}}
\makeatother

\begin{document}
    \frame{
        \titlepage
    }

    \begin{frame}{Plan}
        \tableofcontents
    \end{frame}

    \section{Квадратичні лишки} % Note all sections and subsections are automatically placed in your table of contents

%------------------------------------------------
\begin{frame}
    \frametitle{Символ Лежандра. Повторення}
    \small
    \begin{block}{Означення}
        Нехай $p>2$ -- просте, $a\in\mathbb{Z}$. \textbf{Символ Лежандра}:
        \vspace{-1mm}
        \begin{equation*}
            \legendre{a}{p}=
            \begin{cases}
                \ \ 0, & p\mid a,\\
                \ \ 1, & a\in\mathcal{Q}_p\ \ (\exists x:\ x^2\equiv a \bmod p),\\
                -1, & a\in\overline{\mathcal{Q}}_p.
            \end{cases}
        \end{equation*}
    \end{block}
    \vspace{-2mm}
    \begin{exampleblock}{Закони}
        \begin{enumerate}[itemsep=2pt]
            \item $a\equiv b \pmod p\ \Rightarrow\ \legendre{a}{p}=\legendre{b}{p}$;
            \item $\legendre{ab}{p}=\legendre{a}{p}\legendre{b}{p}$, \ зокрема $\legendre{a^2}{p}=1$ при $p\nmid a$;
            \item критерій Ойлера: $\legendre{a}{p}\equiv a^{\frac{p-1}{2}} \pmod p$;
            \item $\legendre{1}{p}=1,\quad \legendre{-1}{p}=(-1)^{\frac{p-1}{2}},\quad \legendre{2}{p}=(-1)^{\frac{p^2-1}{8}}$;
            \item закон взаємності: для різних непарних простих $p,q$ \
            $\legendre{p}{q}\legendre{q}{p}=(-1)^{\frac{p-1}{2}\cdot\frac{q-1}{2}}$.
        \end{enumerate}
    \end{exampleblock}
\end{frame}
%------------------------------------------------
\begin{frame}
    \frametitle{Символ Якобі. Псевдоквадрати}
    \footnotesize
    \begin{block}{Означення}
        Нехай $n$ -- непарне, $n=p_1^{\alpha_1}\cdots p_k^{\alpha_k}$. \textbf{Символ Якобі}: \
        $\legendre{a}{n}=\legendre{a}{p_1}^{\alpha_1}\cdots\legendre{a}{p_k}^{\alpha_k}$.
    \end{block}
    \vspace{-2mm}
    \begin{exampleblock}{Закони \normalfont(для непарних $m,n$)}
        \begin{columns}[T]
            \begin{column}{0.47\textwidth}
                \begin{enumerate}[itemsep=1pt]
                    \item $a\equiv b \pmod n\Rightarrow\legendre{a}{n}=\legendre{b}{n}$;
                    \item $\legendre{ab}{n}=\legendre{a}{n}\legendre{b}{n}$;
                    \item $\legendre{a}{mn}=\legendre{a}{m}\legendre{a}{n}$;
                \end{enumerate}
            \end{column}
            \begin{column}{0.5\textwidth}
                \begin{enumerate}[start=4,itemsep=1pt]
                    \item $\legendre{-1}{n}=(-1)^{\frac{n-1}{2}}$, \ $\legendre{2}{n}=(-1)^{\frac{n^2-1}{8}}$;
                    \item $\gcd(m,n)=1$: \ $\legendre{m}{n}\legendre{n}{m}=(-1)^{\frac{m-1}{2}\cdot\frac{n-1}{2}}$.
                \end{enumerate}
            \end{column}
        \end{columns}
        Обчислюється за $O(\log^2 n)$ \textbf{без розкладу} $n$ на множники.
    \end{exampleblock}
    \vspace{-2mm}
    \begin{alertblock}{Псевдоквадрати}
        $\legendre{a}{n}=-1\Rightarrow a\notin\QR_n$, \ але \ $\legendre{a}{n}=1\not\Rightarrow a\in\QR_n$.

        \textbf{Означення.} $a\in\mathbb{Z}_n^*$ -- \textbf{псевдоквадрат} за модулем $n$, якщо $\legendre{a}{n}=1$, але $a\notin\QR_n$. Для $n=pq$:
        \vspace{-1mm}
        \begin{equation*}
            \PSQ_n=\left\{a\in\mathbb{Z}_n^*:\ \legendre{a}{p}=\legendre{a}{q}=-1\right\},\qquad |\QR_n|=|\PSQ_n|=\tfrac{\varphi(n)}{4}.
        \end{equation*}
        \vspace{-4mm}
    \end{alertblock}
\end{frame}
%------------------------------------------------
\begin{frame}
    \frametitle{Задачі}
    \begin{block}{}
        \textbf{1.} Oбчисліть
        $\legendre{74}{437}$ та $\legendre{67}{437}$.
        Що можна сказати про те, чи є ці числа квадратичними лишками за модулем $437=19\cdot 23$?
    \end{block}
    \begin{figure}
        \centering
        \includegraphics[width=1\linewidth]{2-crypto-sym/images/cells.png}
    \end{figure}
\end{frame}
%------------------------------------------------
\section{Протокол Діффі-Геллмана}
%------------------------------------------------
\begin{frame}
    \frametitle{Протокол узгодження ключів}
    \centering
    \begin{tikzpicture}[
        every node/.style={font=\small},
        msg/.style={-{Stealth[length=2.5mm]}, line width=0.9pt},
        party/.style={draw, rounded corners=3pt, line width=0.8pt, minimum width=2.4cm, minimum height=0.65cm, font=\normalsize\bfseries, text=white}
      ]
        \node[draw=oc-gray-6, fill=oc-gray-1, rounded corners=3pt, text width=9.6cm, align=center, inner sep=4pt] (pp) at (4.2,1.05)
            {\textbf{Публічні параметри:} просте $p$, $g\in\mathbb{Z}_p^*$ великого порядку $q$ \\ Покажіть, що обчислений секрет співпадає};

        \node[party, fill=alicecol, draw=alicecol] (A) at (0,0) {Аліса};
        \node[party, fill=bobcol, draw=bobcol] (B) at (8.4,0) {Боб};

        \node[text=alicecol] (a1) at (0,-0.85) {$a\inR\mathbb{Z}_q$};
        \node[text=alicecol] (a2) at (0,-1.55) {$A:=g^{a}\bmod p$};
        \node[text=bobcol]   (b1) at (8.4,-0.85) {$b\inR\mathbb{Z}_q$};
        \node[text=bobcol]   (b2) at (8.4,-2.35) {$B:=g^{b}\bmod p$};

        \draw[msg, alicecol] (1.45,-1.55) -- node[above]{$A$} (6.95,-1.55);
        \draw[msg, bobcol]   (6.95,-2.35) -- node[above]{$B$} (1.45,-2.35);

        \node[text=alicecol] (ka) at (0,-3.25) {$K:=B^{a}\bmod p$};
        \node[text=bobcol]   (kb) at (8.4,-3.25) {$K:=A^{b}\bmod p$};

        \node[draw=oc-yellow-8, fill=oc-yellow-1, rounded corners=3pt, inner sep=4pt] (k) at (4.2,-3.25) {$K=g^{ab}\bmod p$};
        \draw[oc-yellow-8, line width=0.8pt, -{Stealth}] (ka) -- (k);
        \draw[oc-yellow-8, line width=0.8pt, -{Stealth}] (kb) -- (k);

        \node[text=oc-gray-7, font=\footnotesize, align=center] at (4.2,-4.15)
            {По каналу можна відслідковувати усе, що передається у відкриту: $p,\ g,\ A,\ B$};
    \end{tikzpicture}
\end{frame}
%------------------------------------------------
% --- Альтернативна версія попереднього слайду ("ghoul edition"). Залиште одну з двох. ---
\renewcommand{\ghoulbackground}{%
  \begin{tikzpicture}[
      x=1cm, y=1cm,
      txt/.style={text=white, font=\small},
      glass/.style={fill=black, fill opacity=0.6, text opacity=1, rounded corners=2pt, inner sep=3pt},
      msg/.style={-{Stealth[length=2.8mm]}, line width=1.1pt, ghoulred}
    ]
    \useasboundingbox (0,0) rectangle (\paperwidth,\paperheight);
    \ghoulbase{Протокол узгодження ключів}
    \node[anchor=north west, txt, font=\footnotesize, glass] at (0.50,8.3)
      {\textcolor{ghoulred}{\textbf{Публічно:}} просте $p$, \ $g\in\mathbb{Z}_p^*$ порядку $q$};

    % Аліса та Боб
    \ghoulcard{(2.0,5.75)}{3.0cm}{6-crypto-practice/Alice.png}
    \coordinate (Acard) at (card.east);
    \ghoulcard{(10.8,5.75)}{3.0cm}{6-crypto-practice/Bob.png}
    \coordinate (Bcard) at (card.west);

    % повідомлення
    \draw[msg] (3.65,6.55) -- node[above=1pt, txt, glass]{$A:=g^{a}\bmod p$} (9.15,6.55);
    \draw[msg] (9.15,5.15) -- node[above=1pt, txt, glass]{$B:=g^{b}\bmod p$} (3.65,5.15);

    % обчислення сторін
    \node[txt, glass] at (2.0,3.95) {$a\inR\mathbb{Z}_q$};
    \node[txt, glass] at (10.8,3.95) {$b\inR\mathbb{Z}_q$};
    \node[txt, glass] (ka) at (2.0,3.3) {$K:=B^{a}\bmod p$};
    \node[txt, glass] (kb) at (10.8,3.3) {$K:=A^{b}\bmod p$};
    \node[draw=ghoulred, fill=ghoulblood, fill opacity=0.9, text opacity=1, line width=1pt, inner sep=5pt,
          text=white, font=\normalsize\bfseries] (k) at (6.4,3.3) {$K=g^{ab}\bmod p$};
    \draw[ghoulred, line width=0.9pt, -{Stealth}] (ka) -- (k);
    \draw[ghoulred, line width=0.9pt, -{Stealth}] (kb) -- (k);

    % зловмисник слухає канал
    \ghoulcard{(6.4,1.3)}{2.2cm}{6-crypto-practice/Adversar.png}
    \node[txt, font=\footnotesize, glass, anchor=east] at (5.15,1.55) {бачить: \ $p,\ g,\ A,\ B$};
    \node[txt, font=\footnotesize, glass, anchor=west] at (7.65,1.55) {не знає: \ $a,\ b$};
    \node[txt, font=\footnotesize, glass, anchor=east, text=ghoulred!40!white] at (5.15,0.95) {мета: \ $g^{ab}$};

    % пасхалка
    \node[anchor=south east, font=\tiny\ttfamily, text=ghoulgray] at (12.6,0.15) {1000 - 7 = 993 ...};
  \end{tikzpicture}%
}
\begin{frame}[plain, ghoul]
\end{frame}
%------------------------------------------------
\begin{frame}
    \frametitle{Безпека DH}
    \small
    Нехай $G=\langle g\rangle$ -- циклічна група порядку $q$, $a,b,c\inR\mathbb{Z}_q$.
    \begin{block}{Задача Діффі-Геллмана (злам схеми DH)}
        Знаючи $(G,\,g,\,g^{a},\,g^{b})$ -- тобто все, що бачить Єва, -- знайти спільний ключ $g^{ab}$.
    \end{block}
    \vspace{-2mm}
    \begin{block}{CDH (Computational Diffie-Hellman)}
        Для будь-якого PPT-алгоритму $\mathcal{A}$: \
        $\Pr\left[\mathcal{A}(g,g^{a},g^{b})=g^{ab}\right]\le\mathsf{negl}(\lambda)$.
    \end{block}
    \vspace{-2mm}
    \begin{block}{DDH (Decisional Diffie-Hellman)}
        Набори $(g,g^{a},g^{b},g^{ab})$ та $(g,g^{a},g^{b},g^{c})$ обчислювально нерозрізненні:
        \vspace{-1mm}
        \begin{equation*}
            \left|\Pr\left[\mathcal{A}(g,g^{a},g^{b},g^{ab})=1\right]-\Pr\left[\mathcal{A}(g,g^{a},g^{b},g^{c})=1\right]\right|\le\mathsf{negl}(\lambda).
        \end{equation*}
    \end{block}
    \vspace{-2mm}
    \begin{alertblock}{}
        \footnotesize
        Розв'язуємо DLP $\Rightarrow$ розв'язуємо CDH $\Rightarrow$ розв'язуємо DDH. Отже, припущення DDH -- найсильніше з трьох.
        У всій групі $\mathbb{Z}_p^*$ DDH \textbf{не виконується}: $\legendre{g^{x}}{p}=(-1)^{x}$ видає парність $x$.
    \end{alertblock}
\end{frame}
%------------------------------------------------
\begin{frame}
    \frametitle{Атаки на DH: людина посередині (MITM)}
    \centering
    \begin{tikzpicture}[
        every node/.style={font=\small},
        msg/.style={-{Stealth[length=2.5mm]}, line width=0.9pt},
        party/.style={draw, rounded corners=3pt, line width=0.8pt, minimum width=2.2cm, minimum height=0.6cm, font=\normalsize\bfseries, text=white}
      ]
        \node[party, fill=alicecol, draw=alicecol] (A) at (0,0) {Аліса};
        \node[party, fill=evecol, draw=evecol] (E) at (4.6,0) {Єва};
        \node[party, fill=bobcol, draw=bobcol] (B) at (9.2,0) {Боб};

        \node[text=alicecol] at (0,-0.8) {$a\inR\mathbb{Z}_q$};
        \node[text=evecol]   at (4.6,-0.8) {$e,\,e'\inR\mathbb{Z}_q$};
        \node[text=bobcol]   at (9.2,-0.8) {$b\inR\mathbb{Z}_q$};

        \draw[msg, alicecol] (0.4,-1.45) -- node[above]{$g^{a}$} (4.2,-1.45);
        \draw[msg, evecol]   (5.0,-1.45) -- node[above]{$g^{e}$} (8.8,-1.45);
        \draw[msg, bobcol]   (8.8,-2.25) -- node[above]{$g^{b}$} (5.0,-2.25);
        \draw[msg, evecol]   (4.2,-2.25) -- node[above]{$g^{e'}$} (0.4,-2.25);

        \draw[oc-gray-5, dashed] (4.6,-1.1) -- (4.6,-2.6);
        \draw[oc-gray-5, dashed] (9.2,-1.1) -- (9.2,-2.6);
        \draw[oc-gray-5, dashed] (0,-1.1) -- (0,-2.6);

        \node[text=alicecol] at (0,-3.0) {$K_A=(g^{e'})^{a}$};
        \node[text=bobcol]   at (9.2,-3.0) {$K_B=(g^{e})^{b}$};
        \node[text=evecol, align=center] at (4.6,-3.15) {$K_A=(g^{a})^{e'}$\\ $K_B=(g^{b})^{e}$};
    \end{tikzpicture}

    \vspace{1mm}
    \begin{alertblock}{}
        \footnotesize
        Аліса і Боб упевнені, що мають спільний ключ, але насправді кожен узгодив ключ з Євою: вона розшифровує і перешифровує весь трафік.
        Причина -- \textbf{відсутність автентифікації}: CDH тут ні до чого, Єва нічого не логарифмує.
        Захист: підписи на повідомленнях протоколу / сертифікати (STS, TLS), статичні ключі ідентичності (X3DH).
    \end{alertblock}
\end{frame}
%------------------------------------------------
% --- Альтернативна версія попереднього слайду ("ghoul edition"). Залиште одну з двох. ---
\renewcommand{\ghoulbackground}{%
  \begin{tikzpicture}[x=1cm, y=1cm, txt/.style={gtxt}, msg/.style={gmsg}]
    \useasboundingbox (0,0) rectangle (\paperwidth,\paperheight);
    \ghoulbase{Атаки на DH: людина посередині}
    \node[anchor=north west, txt, font=\footnotesize, glass] at (0.50,8.3)
      {Зловмисник \textcolor{ghoulred}{\textbf{контролює канал}}: перехоплює та підміняє повідомлення};

    % учасники
    \ghoulcard{(1.85,5.55)}{2.6cm}{6-crypto-practice/Alice.png}
    \ghoulcard{(6.4,5.55)}{2.6cm}{6-crypto-practice/Adversar.png}
    \ghoulcard{(10.95,5.55)}{2.6cm}{6-crypto-practice/Bob.png}

    % повідомлення
    \draw[msg] (3.3,6.15) -- node[above=1pt, txt, glass]{$g^{a}$} (4.95,6.15);
    \draw[msg] (7.85,6.15) -- node[above=1pt, txt, glass]{$g^{e}$} (9.5,6.15);
    \draw[msg] (9.5,4.95) -- node[above=1pt, txt, glass]{$g^{b}$} (7.85,4.95);
    \draw[msg] (4.95,4.95) -- node[above=1pt, txt, glass]{$g^{e'}$} (3.3,4.95);

    % вибір секретів
    \node[txt, glass] at (1.85,3.85) {$a\inR\mathbb{Z}_q$};
    \node[txt, glass] at (6.4,3.85) {$e,\,e'\inR\mathbb{Z}_q$};
    \node[txt, glass] at (10.95,3.85) {$b\inR\mathbb{Z}_q$};

    % ключі
    \node[txt, glass] at (1.85,3.2) {$K_A:=(g^{e'})^{a}$};
    \node[txt, glass] at (10.95,3.2) {$K_B:=(g^{e})^{b}$};
    \node[draw=ghoulred, fill=ghoulblood, fill opacity=0.9, text opacity=1, line width=1pt, inner sep=4pt,
          text=white, font=\small, align=center] at (6.4,3.0)
          {$K_A=(g^{a})^{e'}$\\[1pt] $K_B=(g^{b})^{e}$};

    % пояснення
    \node[txt, font=\footnotesize, glass, text width=11.4cm, align=left, inner sep=5pt] at (6.4,1.05)
      {Аліса і Боб упевнені, що мають спільний ключ, але кожен узгодив ключ зі Зловмисником:
       він розшифровує і перешифровує весь трафік. Причина -- \textcolor{ghoulred!80!white}{\textbf{відсутність автентифікації}}:
       CDH тут ні до чого, нічого не логарифмується.\\
       \textcolor{ghoulred!80!white}{\textbf{Захист:}} підписи на повідомленнях протоколу / сертифікати (STS, TLS), ключі ідентичності (X3DH).};
  \end{tikzpicture}%
}
\begin{frame}[plain, ghoul]
\end{frame}
%------------------------------------------------
\begin{frame}
    \frametitle{X3DH (Extended Triple Diffie-Hellman, Signal)}
    \footnotesize
    \begin{columns}[T]
        \begin{column}{0.46\textwidth}
            Асинхронне узгодження: Боб офлайн, його ключі -- на сервері.
            \begin{block}{Ключі}
                $IK$ -- ключ ідентичності (довгостроковий);\\
                $SPK_B$ -- prekey Боба + $\mathsf{Sig}_{IK_B}(SPK_B)$;\\
                $OPK_B$ -- одноразовий prekey (опційно);\\
                $EK_A$ -- ефемерний ключ Аліси.
            \end{block}
        \end{column}
        \begin{column}{0.52\textwidth}
            \centering
            \begin{tikzpicture}[
                key/.style={draw, rounded corners=2pt, line width=0.8pt, minimum width=1.25cm, minimum height=0.5cm, font=\footnotesize\bfseries, fill=oc-gray-0},
                dh/.style={line width=0.9pt},
                lab/.style={font=\scriptsize, fill=oc-gray-0, inner sep=1pt}
              ]
                \node[font=\small\bfseries, text=alicecol] at (0,0.6) {Аліса};
                \node[font=\small\bfseries, text=bobcol] at (4.0,0.6) {Боб};
                \node[key, draw=alicecol, text=alicecol] (IKA) at (0,0) {$IK_A$};
                \node[key, draw=alicecol, text=alicecol] (EKA) at (0,-2.0) {$EK_A$};
                \node[key, draw=bobcol, text=bobcol] (IKB) at (4.0,0) {$IK_B$};
                \node[key, draw=bobcol, text=bobcol] (SPKB) at (4.0,-1.2) {$SPK_B$};
                \node[key, draw=bobcol, text=bobcol, dashed] (OPKB) at (4.0,-2.4) {$OPK_B$};

                \draw[dh, oc-indigo-6] (IKA.east) -- node[lab, pos=0.2]{$DH_1$} (SPKB.west);
                \draw[dh, oc-indigo-6] (EKA.east) -- node[lab, pos=0.8]{$DH_2$} (IKB.west);
                \draw[dh, oc-orange-7] (EKA.east) -- node[lab, pos=0.75]{$DH_3$} (SPKB.west);
                \draw[dh, oc-orange-7, dashed] (EKA.east) -- node[lab, pos=0.55]{$DH_4$} (OPKB.west);
            \end{tikzpicture}
        \end{column}
    \end{columns}

    \vspace{-1mm}
    \begin{exampleblock}{}
        $SK=\mathsf{KDF}\left(DH_1\,\|\,DH_2\,\|\,DH_3\,\|\,DH_4\right)$
        \begin{itemize}[itemsep=1pt]
            \item \textcolor{oc-indigo-7}{$DH_1, DH_2$} -- взаємна автентифікація (захист від MITM);
            \item \textcolor{oc-orange-8}{$DH_3, DH_4$} -- пряма секретність (forward secrecy);
            \item Аліса надсилає $IK_A$, $EK_A$, ідентифікатори використаних prekey та перше повідомлення під $SK$ (AEAD).
        \end{itemize}
    \end{exampleblock}
    {\scriptsize Сьогодні Signal використовує PQXDH: до $DH_1,\ldots,DH_4$ додається спільний секрет постквантового KEM (ML-KEM).}
\end{frame}
%------------------------------------------------
\begin{frame}
    \frametitle{Задачі}
    \begin{block}{}
        \textbf{2.} Нехай $p$ -- просте, $g\in\mathbb{Z}_p^*$ -- генератор. Покажіть, що за значенням $y=g^{x}\bmod p$ Єва завжди може сказати, парним чи непарним є невідоме їй $x$.

        \textbf{3.} Переформулюйте схему Діффі-Геллмана в термінах групи $(\mathbb{Z}_n,+)$. Чи буде стійкою така схема?
    \end{block}
    \begin{figure}
        \centering
        \includegraphics[width=1\linewidth]{2-crypto-sym/images/cells.png}
    \end{figure}
\end{frame}
%------------------------------------------------
\section{Протокол шифрування RSA}
%------------------------------------------------
\begin{frame}
    \frametitle{Визначення параметрів та схема}
    \small
    \begin{block}{Генерування ключів}
        Обрати великі прості $p\neq q$, обчислити $n=pq$, $\varphi(n)=(p-1)(q-1)$.
        Обрати $e$: $\gcd(e,\varphi(n))=1$, обчислити $d=e^{-1}\bmod\varphi(n)$.\\
        Відкритий ключ -- $(n,e)$, секретний ключ -- $d$ (а також $p,q$).
    \end{block}
    \vspace{-1mm}
    \centering
    \begin{tikzpicture}[
        every node/.style={font=\small},
        msg/.style={-{Stealth[length=2.5mm]}, line width=0.9pt},
        party/.style={draw, rounded corners=3pt, line width=0.8pt, minimum width=2.4cm, minimum height=0.6cm, font=\normalsize\bfseries, text=white}
      ]
        \node[party, fill=bobcol, draw=bobcol] (B) at (0,0) {Боб};
        \node[party, fill=alicecol, draw=alicecol] (A) at (8.4,0) {Аліса};
        \node[text=alicecol] at (8.4,-0.8) {$(n,e),\ d$};
        \draw[msg, alicecol] (7.5,-0.8) -- node[above]{$pk=(n,e)$} (1.2,-0.8);
        \node[text=bobcol] at (0,-0.8) {$m\in\mathbb{Z}_n$};
        \node[text=bobcol] at (0,-1.6) {$c:=m^{e}\bmod n$};
        \draw[msg, bobcol] (1.3,-1.6) -- node[above]{$c$} (7.1,-1.6);
        \node[text=alicecol] at (8.4,-1.6) {$m:=c^{d}\bmod n$};
    \end{tikzpicture}
    \vspace{-1mm}
    \begin{alertblock}{}
        \raggedright
        Коректність при $\gcd(m,n)=1$ (загальний випадок -- задача 6): \\
        $ed=1+k\varphi(n)\ \Rightarrow\ c^{d}=m^{ed}=m\cdot\left(m^{\varphi(n)}\right)^{k}\equiv m \pmod n$.\\
        Підпис: $s=m^{d}\bmod n$, перевірка: $s^{e}\overset{?}{\equiv}m\pmod n$.
    \end{alertblock}
\end{frame}
%------------------------------------------------
\begin{frame}
    \frametitle{Задачі}
    \begin{block}{}
        \textbf{4.} Переформулюйте криптосистему RSA для випадку простого модуля.
        \begin{enumerate}
            \item[a)] покажіть, що така система не є стійкою;
            \item[b)] покажіть, що якщо обрати $p=q$, то опис криптосистеми стане некоректним, а система стане нестійкою.
        \end{enumerate}
        \textbf{5.} Покажіть, як за відомим $n$ та значенням функції Ойлера $\varphi(n)$ обчислити значення $p$ та $q$.
    \end{block}
    \begin{figure}
        \centering
        \includegraphics[width=1\linewidth]{2-crypto-sym/images/cells_small.png}
    \end{figure}
\end{frame}
%------------------------------------------------
\begin{frame}
    \frametitle{Задачі}
    \begin{block}{}
        \textbf{6.} Доведіть, що процедура розшифрування RSA коректно розшифровує відкритий текст $m$ навіть якщо $\gcd(m,n)>1$. Оцініть ймовірність того, що випадково обране повідомлення буде невзаємнопростим із модулем $n$.

        \textbf{7.} Яка складність перевіряння цифрового підпису в схемі RSA для значення $e=2^{t}+1$? Наведіть числові відповіді для $e=3,5,17,65537$. В якості базової операції використовується множення двох чисел за модулем.
    \end{block}
    \begin{figure}
        \centering
        \includegraphics[width=1\linewidth]{2-crypto-sym/images/cells_small.png}
    \end{figure}
\end{frame}
%------------------------------------------------
\begin{frame}
    \frametitle{Задачі}
    \begin{block}{}
        \textbf{8.} Покажіть, що при $e=\frac{\varphi(n)}{2}+1$ функція RSA перетворюється на тотожну.

        \textbf{9.} Доведіть, що функція RSA зберігає символ Якобі: $\legendre{m^{e}}{n}=\legendre{m}{n}$.

        \textbf{10.} Нехай $c=m^{e}\bmod n$ -- шифротекст RSA. Покажіть, що для довільного $u\in\mathbb{Z}_n$ число $c':=c\cdot u^{e}\bmod n$ є коректним шифротекстом під повідомленням $m':=m\cdot u\bmod n$.
    \end{block}
    \begin{figure}
        \centering
        \includegraphics[width=1\linewidth]{2-crypto-sym/images/cells_small.png}
    \end{figure}
\end{frame}
%------------------------------------------------
\section{Схема Горнера}
%------------------------------------------------
\begin{frame}
    \frametitle{Схема}
    \small
    \textbf{Вхід:} $a,m\in\mathbb{N}\ (a<m)$, \ $n$, \ $|n|=k$ біт. \quad \textbf{Обчислюємо:} $a^{n}\bmod m$.
    \begin{block}{Ідея}
        \vspace{-4mm}
        \begin{gather*}
            n=(n_{k-1}n_{k-2}\ldots n_1n_0)_2=\textstyle\sum_{i=0}^{k-1}n_i2^{i}\\
            a^{n}=a^{\sum_{i=0}^{k-1}n_i2^{i}}=\Big(\ldots\big(\big(a^{n_{k-1}}\big)^{2}\cdot a^{n_{k-2}}\big)^{2}\cdot a^{n_{k-3}}\ldots\Big)^{2}\cdot a^{n_0}
        \end{gather*}
        \vspace{-5mm}

        {\footnotesize Після кожного піднесення до квадрату та множення беремо число за $\bmod\ m$.}
    \end{block}
    \vspace{-2mm}
    \begin{columns}[T]
        \begin{column}{0.52\textwidth}
            \begin{exampleblock}{Ітеративно}
                \begin{enumerate}
                    \item $b:=1$
                    \item \textbf{for} $i=k-1$ \textbf{downto} $0$:\\
                          \quad $b:=b^{2}\bmod m$\\
                          \quad \textbf{if} $n_i=1$: \ $b:=b\cdot a\bmod m$
                    \item \textbf{return} $b$
                \end{enumerate}
            \end{exampleblock}
        \end{column}
        \begin{column}{0.44\textwidth}
            \begin{alertblock}{Складність}
                $k$ піднесень до квадрату $+$ $\mathrm{wt}(n)\le k$ множень: \ $\le 2k$ модулярних множень замість $n-1$.
            \end{alertblock}
        \end{column}
    \end{columns}
\end{frame}
%------------------------------------------------
\begin{frame}
    \frametitle{Приклад: $23^{143}\bmod 41$}
    \small
    $143=$
    \includegraphics[width=1\linewidth]{2-crypto-sym/images/cells.png}
\end{frame}
%------------------------------------------------
\begin{frame}
    \frametitle{Задачі}
    \begin{block}{}
        \textbf{11.} Побудуйте спільний секрет для абонентів A та B у схемі Діффі-Геллмана, якщо $p=11,\ g=2,\ x_A=5,\ x_B=7$.
    \end{block}
    % Відповіді: 11) y_A = 2^5 = 10, y_B = 2^7 = 7 (mod 11), K = 2^35 = 10 (mod 11).
    \begin{figure}
        \centering
        \includegraphics[width=1\linewidth]{2-crypto-sym/images/cells.png}
    \end{figure}
\end{frame}

\begin{frame}{Задачі}
    \begin{block}{}
        \textbf{12.} Розшифруйте повідомлення $C=815$, якщо відомо, що використовується схема RSA, $p=37,\ q=43$, а відкритий ключ $e=17$.
    % Відповідь: n = 1591, phi(n) = 1512, d = 17^{-1} mod 1512 = 89 = (1011001)_2, m = 815^89 mod 1591 = 1000.
    \end{block}
        \begin{figure}
        \centering
        \includegraphics[width=1\linewidth]{2-crypto-sym/images/cells.png}
    \end{figure}
\end{frame}

    \begin{frame}[plain, standout]
    \begin{figure}
        \centering
        \includegraphics[width=0.9\linewidth]{Cryptocats.jpg}
        \label{fig:placeholder}
    \end{figure}
    \centering
    \LARGE
    \textbf{Питання?} \\

    \vspace{0.2cm} \Huge \ding{170} \large \\

  \end{frame}
\end{document}
